\section{Conclusion}
\label{sec:conclusion}

Deleting multi-agent roles largely eliminates the intended division of
labor for every model we tested on Omni-MATH, without lowering first-pass
accuracy. Besides dividing work, roles allocate decision rights. A
role-structured Reviewer controls when the team submits to the external
evaluator. In exploratory comparisons with equal access to feedback, the
role-free team recovers from a failed first answer more often in every
tested setting. A post hoc comparison at equal numbers of verdicts removes
this advantage, and so does forcing gemma-4's role-structured team to
submit in every cycle. Roles do not just assign responsibilities; they also
govern decisions and information flow, which shape how a team recovers
from failure. Building an effective multi-agent organization therefore goes
beyond assigning roles correctly. It requires designing decision rights and
information pathways that support the desired system behavior.
